\subsection{
    Mecanismos de Defesa
} \label{sec:mecanismos_defesa}

Os mecanismos de defesa complementam a taxonomia de vulnerabilidades ao fornecer técnicas, ferramentas e controles concretos para prevenir, detectar ou mitigar ataques. Esta seção organiza as defesas por categoria de ataque, apresenta as ferramentas de análise estática e dinâmica, e discute as soluções de proteção em tempo de execução (\textit{runtime}).

\subsubsection{Defesas por categoria}

Cada categoria de vulnerabilidade descrita nas seções anteriores possui mecanismos de defesa específicos, com diferentes níveis de efetividade e limitações. A Tabela \ref{tab:tecnicas_defesa} consolida a relação entre técnicas de defesa, ataques mitigados e suas restrições, com base em trabalhos de classificação como os de \cite{halfond:2006}, \cite{gupta:2015}, \cite{barth:2008} e \cite{owasp:2021}.

\begin{table}[htbp]
    \caption{Técnicas de defesa por categoria de ataque e limitações}
    \label{tab:tecnicas_defesa}
    \centering
\small
\begin{tabularx}
    {\textwidth}
    {
        @{}
        >{\hsize=0.75\hsize}X 
        >{\hsize=0.75\hsize}X 
        >{\hsize=1.5\hsize}X
        @{}
    }
    \toprule
        \textbf{Técnica} 
        & \textbf{Ataque} 
        & \textbf{Limitações} \\ 
    \midrule
    \csvreader
        [head to column names, late after line=\\ \addlinespace]
        {tabelas/dados/tabela_tecnicas_defesa.csv}
        {}
        {
            \tecnica 
            & \ataque 
            & \limitacoes
        }
    \bottomrule
    \multicolumn
        {3}
        {l}
        {
            \footnotesize 
            \textbf{Fonte:}
            \cite{Halfond2006ClassificationSQLInjectionAttacksCountermeasures};
            \cite{Heiderich2013MXSSAttacksAttackingWellSecuredWebApplications};
            \cite{Gupta2015CrossSiteScriptingAttacksDefenseMechanisms};
            \cite{Barth2022SystematicApproachWebProtectionRuntimeTools};
            \cite{None2008RobustDefensesCrossSiteRequestForgery}.
        }
\end{tabularx}

\end{table}

\subsubsection{Ferramentas de análise: SAST, DAST, IAST}

As ferramentas de análise de segurança identificam vulnerabilidades ao longo do ciclo de vida da aplicação. O SAST (\textit{Static Application Security Testing}) analisa código-fonte ou \textit{bytecode} sem execução, sendo adequado para integração em pipelines CI/CD \cite{alazmi:2022}. O DAST (\textit{Dynamic Application Security Testing}) realiza testes em tempo de execução, simulando ataques à aplicação, com avanços como a automação de autenticação via OWASP ZAP (Sacramento et al., 2024). Já o IAST (\textit{Interactive Application Security Testing}) combina abordagens estática e dinâmica por meio de instrumentação durante testes funcionais, aumentando a precisão e reduzindo falsos positivos \cite{seth:2025}. Evidências recentes reforçam a importância da combinação dessas técnicas para uma cobertura mais eficaz \cite{ieee:2024}.

\begin{table}[htbp]
    \caption{Comparativo SAST, DAST, IAST}
    \label{tab:comparativo_sast_dast_iast}
    \centering
    \small
    \begin{tabularx}
        {\textwidth}
        {@{}XXXX@{}}
        \toprule
            \textbf{Critério}
            & \textbf{SAST}
            & \textbf{DAST}
            & \textbf{IAST} \\ 
        \midrule
        \csvreader
            [head to column names, late after line=\\ ]
            {tabelas/dados/tabela_comparativo_sast_dast_iast.csv}
            {}
            {
                \criterio
                & \sast
                & \dast
                & \iast
            }
        \bottomrule
        \multicolumn
            {4}
            {l}
            {
                \footnotesize
                \textbf{Fonte:}
                \cite{Alazmi2022WebScannerReview};
                \cite{Sacramento2024AutenticacaoZAP};
                \cite{DeJesus2023SecurityTestingWebApps};
                \cite{Cotroneo2025IASTvsRASP}.
            }
    \end{tabularx}

\end{table}

\subsubsection{Defesas em \textit{runtime}: WAF, WAAP, RASP}

Defesas em runtime atuam em produção para proteger a aplicação em execução. O WAF (\textit{Web Application Firewall}) inspeciona o tráfego HTTP e bloqueia requisições maliciosas com base em regras e assinaturas, embora enfrente limitações semelhantes às de ferramentas \textit{black-box} \cite{Doupe_Cova_Vigna_2010}. O WAAP (\textit{Web Application and API Protection}) amplia essa abordagem ao incluir proteção de APIs e mitigação de ataques modernos. Já o RASP (\textit{Runtime Application Self-Protection}) atua dentro da aplicação, utilizando o contexto de execução para aumentar a precisão, reduzir falsos positivos e mitigar ataques \textit{zero-day} \cite{sureda:2022}.

\subsubsection{Comparativo WAF × RASP}

\cite{sureda:2022} realizaram uma análise sistemática da efetividade de ferramentas de proteção em \textit{runtime}, comparando WAF e RASP em diferentes categorias de ataque. Os resultados indicam que o RASP supera o WAF em precisão por categoria de ataque, em parte devido ao conhecimento contextual da aplicação. A Tabela  \ref{tab:comparativo_waf_rasp} sintetiza o comparativo empírico com base nesse estudo.

\begin{table}[htbp]
    \caption{Comparativo empírico WAF × RASP}
    \label{tab:comparativo_waf_rasp}
    \centering
\small
\begin{tabularx}
    {\textwidth}
    {@{}XXXX@{}}
    \toprule
        \textbf{Categoria \newline de ataque} 
        & \textbf{WAF \newline (efetividade)} 
        & \textbf{RASP \newline (efetividade)} 
        & \textbf{Observação}\\ 
    \midrule
    \csvreader
        [head to column names, late after line=\\ \addlinespace]
        {tabelas/dados/tabela_comparativo_waf_rasp.csv}
        {}
        {
            \categoria
            & \waf
            & \rasp
            & \observacao
        }
    \bottomrule
    \multicolumn
        {4}
        {l}
        {
            \textbf{Fonte:}
            \footnotesize
            \cite{Alazmi2022WebScannerReview};
            \cite{Sacramento2024AutenticacaoZAP};
            \cite{DeJesus2023SecurityTestingWebApps};
            \cite{Cotroneo2025IASTvsRASP}.
        }
\end{tabularx}

\end{table}

Os autores recomendam o uso complementar de WAF e RASP: o WAF na borda para filtrar tráfego conhecido e mitigar DDoS, e o RASP para proteção contextual contra explorações que ultrapassem a camada perimetral.
